\subsection{Objectif comme objet formel de premier rang}
Chaque tâche est encodée dans un objet versionné
\[
G=(g,\mathcal{D},\mathcal{C},\mathcal{R},\pi,h,\rho,\ell,\mathcal{A},\sigma),
\]
où \(g\) est la cible, \(\mathcal{D}\) le domaine, \(\mathcal{C}\) les contraintes, \(\mathcal{R}\) les ressources, \(\pi\) la politique de priorité, \(h\) l'horizon, \(\rho\) le risque acceptable, \(\ell\) le niveau de preuve requis, \(\mathcal{A}\) les critères d'acceptation et \(\sigma\) la politique d'arrêt. Les préférences ordonnent le travail, mais ne neutralisent jamais une précondition, une restriction de capacité ou une porte de sûreté exigée par le contrat. Le synthétiseur doit chercher dans l'espace des compositions compatibles avec \(G\), et tracer les raisons pour lesquelles un candidat est admis, écarté ou conservé comme non résolu.
